The prescribed-image qualification in Theorem~\ref{AL03} is
essential.  In its target, replace $x$ by $T^{[a]}x$ with
$a>0$.  Then
\[
 \gamma_2\bigl(d(T^{[a]}x)\bigr)
   =\binom{2a}{a}T^{[2a]}\vartheta
   =d\left(\frac{\binom{2a}{a}}{2}T^{[2a]}u\right).
 \tag{AL.19}
\]
The coefficient is integral since
$\binom{2a}{a}=2\binom{2a-1}{a-1}$.
Thus this example does not obstruct changing the original disk,
and it does not obstruct choosing images independently when
that original generator is absent from a delayed subalgebra.
The following all-arity calculation makes this comparison
precise without division in the coefficient ring.

\begin{proposition}[A-L04: Proved; explicit extensions for coefficient-scaled odd disks]
\label{AL04}
Let $R$ be a strict PD differential algebra over any commutative
ring, with arbitrary torsion.  Let $z$ be a closed even element
of form degree zero in its distinguished PD ideal, and let
$x$ be an odd element of that ideal, of form degree $2r-1$
with $r\geq1$.  Write $y=dx$.  For every $n\geq1$ define
\[
 Y_n=\gamma_n(zy),\qquad
 X_n=(n-1)!\,\gamma_n(z)\,x\,\gamma_{n-1}(y).
 \tag{AL.20}
\]
Then $dX_n=Y_n$.  For every positive integer $p$ and every
$n\geq1$ one has the exact identity
\[
 \frac{(pn)!}{(n!)^p}\,X_{pn}
       =X_nY_n^{p-1}.
 \tag{AL.21}
\]
The displayed factorial ratio is an integer.
Consequently, over a $\mathbb Z_{(p)}$-algebra, the $p$-local
one-generator closure of SA08, C18, with original generator
mapping to $zx$, has the explicit
strict PD differential extension
$x_i\mapsto X_{p^i}$ and $v_i\mapsto0$ for all $i\geq0$.
For homogeneous weights $\operatorname{wt}(z)=a$ and
$\operatorname{wt}(x)=b$, these assignments have weights
$p^i(a+b)$ and the required form degrees.
\end{proposition}

\begin{proof}
Every expression in \textup{(AL.20)} is integral, and $dz=dy=0$.
The differential rule and the PD product identity give
\[
 dX_n=(n-1)!\gamma_n(z)y\gamma_{n-1}(y)
     =n!\gamma_n(z)\gamma_n(y)
     =z^n\gamma_n(y)=\gamma_n(zy).
 \tag{AL.22}
\]
For \textup{(AL.21)}, the two PD multiplication calculations are
\[
 \gamma_n(z)z^{n(p-1)}
       =\frac{(pn)!}{n!}\gamma_{pn}(z),\qquad
 \gamma_{n-1}(y)\gamma_n(y)^{p-1}
       =\frac{(pn-1)!}{(n-1)!(n!)^{p-1}}
                         \gamma_{pn-1}(y).
 \tag{AL.23}
\]
Both coefficients count ordered partitions with the indicated
block sizes, and hence are integers.  Multiplying these
identities and the factor $(n-1)!$ in $X_n$ gives
\[
 X_nY_n^{p-1}
  =\frac{(pn)!(pn-1)!}{(n!)^p}
            \gamma_{pn}(z)x\gamma_{pn-1}(y)
  =\frac{(pn)!}{(n!)^p}X_{pn}.
\]
No integer has been cancelled in $R$, so the proof applies
with $2$-torsion, $3$-torsion and without flatness.
The closure in C18 has equations
$dx_i=\gamma_{p^i}(y_0)$ and
$dv_i=c_i x_{i+1}-x_i\gamma_{p^i}(y_0)^{p-1}$, with
$c_i=(p^{i+1})!/(p^i!)^p$ and $y_0=d(zx)=zy$.
Equations \textup{(AL.22)--(AL.23)} verify these equations
for the proposed images.  Its underlying graded PD algebra
is freely generated by the specified even and odd generators
over $\Gamma(y_0)$, so the assignments extend by the explicit
free PD construction of Proposition~\ref{B002}.
The generator checks prove that the full extension is a
chain map.  The two degree calculations are
$|X_n|=(2r-1)+(n-1)2r=2rn-1$ and
$\operatorname{wt}(X_n)=na+b+(n-1)b=n(a+b)$.
The zero images of the $v_i$ respect their degrees.
\end{proof}

\begin{theorem}[A-L05: Proved; an integral map from the full scaled bar closure]
\label{AL05}
Let $k$ be any commutative ring and $r\geq1$.  Let $M$ be a
free $k$-module and let $(R,J,\gamma)$ be a strict PD
cochain $k$-algebra.  Specify a closed element
$z\in J^0$ and a $k$-linear map
$\alpha:M\to J^{2r-1}$.  There is a natural strict PD
cochain algebra map
\[
 \mathcal I_{z,\alpha}:U_r(M)\longrightarrow R
 \tag{AL.24}
\]
from the integral one-sided bar closure of SA08, D-029.
On its original disk it sends $x(a)$ to $z\alpha(a)$ and
$y(a)$ to $z\,d\alpha(a)$.
No target exactness, division in $k$, or completion is required.

More explicitly put $B=\Gamma_k(M)$, with divided-power arity,
and let $\phi:B\to R$ be the graded PD map given by
$a\mapsto d\alpha(a)$.  Let $\delta:B\to R$ be the
$\phi$-derivation of form degree minus one given by
$\delta(a)=\alpha(a)$ and
$\delta\gamma_j(a)=\gamma_{j-1}(d\alpha(a))\alpha(a)$.
For a homogeneous bar word
$b_0[b_1|\cdots|b_n]$, put
\[
 a_i=\operatorname{arity}(b_i),\quad
 m=\sum_{i=0}^n a_i,\quad
 A_j=a_j+\cdots+a_n\ (1\leq j\leq n),\quad
 D=\prod_{j=1}^n A_j .
\]
Here $a_0\geq0$, $a_i>0$ for $i>0$, and $D=1$ for $n=0$.
Then
\[
 \mathcal I_{z,\alpha}\bigl(b_0[b_1|\cdots|b_n]\bigr)
  =\frac{m!}{D}\,\gamma_m(z)\,
       \phi(b_0)\delta(b_1)\cdots\delta(b_n).
 \tag{AL.25}
\]
The coefficient $m!/D$ is an integer.
The map is natural for linear maps of $M$ compatible with
$\alpha$; in particular it is simplicial when $M$ is
degreewise free simplicial, $R$ is simplicial, and $z,\alpha$
are simplicially compatible.
If $\operatorname{wt}(z)=a$ and the image of $\alpha$ has
weight $b$, giving the source arity-one disk weight $a+b$
makes \textup{(AL.24)} weight preserving.
\end{theorem}

\begin{proof}
We use the explicit bar object, so the relevant structures
and their gradings are fixed.  Its underlying module is
$\bigoplus_{n\geq0}B\otimes\overline B^{\otimes n}$.
A word of total arity $m$ and bar length $n$ has form
degree $2rm-n$.  Its differential is
\[
 d\bigl(b_0[b_1|\cdots|b_n]\bigr)
 =b_0b_1[b_2|\cdots|b_n]
 +\sum_{i=1}^{n-1}(-1)^i
       b_0[b_1|\cdots|b_i b_{i+1}|\cdots|b_n].
 \tag{AL.26}
\]
Multiplication multiplies $b_0$'s and sums signed shuffles
of the positive bar entries.  These are the formulas of
D-029; the normalized-bar PD structure is also described in
\cite[Theorem 7.5, pp.~250--251; \S8, pp.~252--253]{B_Richter2009}.
Each bar entry has odd regraded degree, which
accounts for every sign in \textup{(AL.26)}.
The divided powers are the integral shuffle divided
powers.  On a positive even-length word their integrality
comes from partitioning the shuffles of its repeated
blocks into free orbits of the block permutation group;
on length-zero words they are those of $B$.  The sum
formula treats every even sum.  The usual cancellation of
successive merges verifies $d^2=0$; the signed shuffle
identities give the multiplication and derivation
identities.  The block argument, followed by rational
verification over the free integral basis and
specialization, proves the PD differential identities.
Thus the explicit bar construction used here has the
stated strict PD structure over every $k$.

The derivation $\delta$ exists by the displayed formula on
divided-power monomials: for a product of divided powers,
sum the expressions obtained by replacing one
$\gamma_j(a)$ by
$\gamma_{j-1}(d\alpha(a))\alpha(a)$.
The other factors are even, so these terms have no
additional signs.  This is the ordinary universal PD
derivation, evaluated on $\alpha$.  The product and sum
identities for divided powers verify that the expression
is basis independent; alternatively construct it in the
free strict PD algebra on symbols $x(a),y(a)$ and use
Proposition~\ref{B002}.
Its defining formulas imply, for homogeneous arity $a$,
\[
 d\phi(b)=0,\qquad d\delta(b)=a\,\phi(b),\qquad
 \delta(bc)=\delta(b)\phi(c)+\phi(b)\delta(c).
 \tag{AL.27}
\]
The middle identity follows on a divided-power monomial
by summing its exponents; the first and last follow
directly from the definition.  They hold over every
coefficient ring.

The suffix sums $A_1>A_2>\cdots>A_n>0$ are distinct
integers at most $m$.  Their product therefore divides
$m!$.  Formula \textup{(AL.25)} is integral.
Its output form degree is
$2ra_0+\sum_{i=1}^n(2ra_i-1)=2rm-n$, and its output
weight is $ma+mb=m(a+b)$.  It therefore has exactly the
required gradings.

We prove its structural identities over the universal
torsion-free integral PD algebra
\[
 \Gamma_{\mathbb Z}(z)\otimes
 \Gamma_{\mathbb Z}(M_y)\otimes
 \Lambda_{\mathbb Z}(M_x),\qquad dz=dy=0,\quad dx=y .
 \tag{AL.28}
\]
Each formula involves finitely many basis elements.
Over its rationalization the coefficient in
\textup{(AL.25)} is $z^m/D$.  Once an identity has been
verified there, injectivity into the rationalization and
the integral formula prove it over \textup{(AL.28)}.
Base change to $k$ and the PD differential map sending
the universal symbols to the prescribed $z,\alpha,d\alpha$
then prove the identity in the given $R$.

For the chain identity fix a word as above.  Differentiating
its image and using \textup{(AL.27)}, the coefficient of
the product with $\delta(b_j)$ replaced by $\phi(b_j)$ is
\[
          (-1)^{j-1}\frac{a_j}{D}.
 \tag{AL.29}
\]
The common factor $z^m\phi(b_0)$ and all other factors are
understood.  In the first term of \textup{(AL.26)}, removing
$b_1$ removes the denominator factor $A_1$.
Merging $b_i,b_{i+1}$ removes the denominator factor
$A_{i+1}$.  Applying the derivation formula in
\textup{(AL.27)}, the coefficient of the same product on
the other side is
$(A_1-A_2)/D$ for $j=1$,
$(-1)^{j-1}(A_j-A_{j+1})/D$ for $1<j<n$,
and $(-1)^{n-1}A_n/D$ for $j=n$.
Set $A_{n+1}=0$; all cases equal
\textup{(AL.29)} because $A_j-A_{j+1}=a_j$.
For $n=1$ only the first term occurs and gives
$A_1/D=a_1/D$.  For $n=0$ both differentials are zero.
This proves the chain identity in every bar length.

For multiplicativity let $a=(a_1,\ldots,a_n)$ and
$b=(b_1,\ldots,b_\ell)$ be the positive arity lists
of two words, and let $D(a),D(b)$ denote their products
of suffix sums.  The required scalar identity is
\[
 \sum_{\sigma\in\operatorname{Sh}(n,\ell)}
       \frac{1}{D(\sigma(a,b))}
            =\frac{1}{D(a)D(b)} .
 \tag{AL.30}
\]
Here shuffle positions are labelled, even when some
arities agree.  To prove it, induct on $n+\ell$.
An empty list gives the identity immediately.
For two nonempty lists, set $A=\sum a_i$, $B=\sum b_i$,
and partition shuffles by whether their first entry
comes from the first or second list.  That first entry
contributes the suffix factor $A+B$.  The induction
hypothesis gives
\[
 \frac{1}{A+B}
 \left(\frac{1}{D(a_{\rm tail})D(b)}
       +\frac{1}{D(a)D(b_{\rm tail})}\right)
 =\frac{A+B}{(A+B)D(a)D(b)} .
\]
This is \textup{(AL.30)}.  Each $\delta(b_i)$ is odd.
The shuffle sign in the source is cancelled by the
same sign when its image is reordered into the two
original lists.  Thus \textup{(AL.30)}, the common power
$z^{m+m'}$, and multiplication of the coefficient
factors prove multiplicativity.  The empty word of
arity zero maps to $1$.  A word of positive arity $m$
maps into $J$, since its formula has the factor
$\gamma_m(z)\in J$.  Hence the distinguished ideals
are also preserved.

Over the rational universal algebra, any unital algebra
map preserves $v^n/n!$ on even elements.  The map
therefore preserves divided powers there.
Both integral PD operations have already been specified,
so injectivity proves compatibility integrally, and
specialization proves it in $R$, including all torsion.
This completes the proof that \textup{(AL.24)} is a
strict PD cochain map.

For an arity-one coefficient $a\in M$, formula
\textup{(AL.25)} gives $z\,d\alpha(a)$ on $a[]$ and
$z\alpha(a)$ on $1[a]$.  Hence it restricts to the
announced scaled disk map.
The formulas for $\phi,\delta$ commute with every linear
map of parameters compatible with $\alpha$; the
suffix factors depend only on arities, which those
maps preserve.  This proves naturality as an equality
of maps, and applying it to each face and degeneracy
proves the simplicial assertion.  The same formulas
also commute with every strict PD differential target
map $R\to R'$: such a map preserves each factor in
\textup{(AL.25)}.  This target naturality applies even
when $\alpha(a)$ is a nonlinear expression in
previously chosen target generators.
\end{proof}

This construction also supplies exact relative compatibility
for boundary maps which already have the displayed form:
for a linear parameter map $i:M_0\to M_1$ and
$\alpha_0=\alpha_1i$, one has
$\mathcal I_{z,\alpha_1}U_r(i)=\mathcal I_{z,\alpha_0}$.
It does not assert such a formula for an arbitrary
previously prescribed boundary map from $U_r(M_0)$.
Likewise, varying $z$ through coefficients of a tail
does not make \textup{(AL.25)} linear in $z$.
Neither the required full coefficient-tail factorization
nor the completion data of C34 have been constructed
by this theorem.

When $R$ has a specified PD differential coefficient map
$S\to R$, the $k$-algebra map in \textup{(AL.24)} has the
coefficient-preserving extension
$S\otimes_kU_r(M)\to R$, $s\otimes v\mapsto s\mathcal I_{z,\alpha}(v)$.
The tensor product has the PD structure extending those
of its two factors; the PD sum and scalar formulas show
directly that this extension preserves it.

\begin{proposition}[A-L06: Proved; independent coefficient parameters and a balancing defect]
\label{AL06}
Let $z\in J^0$ be closed as in Theorem~\ref{AL05}.
Let $L=\bigoplus_{j\in J_0}M_j$, where each $M_j$ is
free over $k$, choose integers $q_j\geq1$, and choose
linear maps $\alpha_j:M_j\to J^{2r-1}$.
There is a natural integral strict PD cochain map
$U_r(L)\to R$ whose disk generator on $M_j$ maps to
$\gamma_{q_j}(z)\alpha_j$.

This map is natural in parameter maps preserving the
$q_j$-grading and in strict PD target maps.
It need not respect an additional $S$-module action
on the collection of coefficient-indexed parameters:
the explicit defect below already occurs with two
odd parameters and coefficients $T,T^{[2]}$.
\end{proposition}

\begin{proof}
Give $M_j$ the additional Euler weight $q_j$.
On $B=\Gamma_k(L)$ define
$\phi(a)=\gamma_{q_j}(z)d\alpha_j(a)$ and the
$\phi$-derivation
$\delta(a)=q_j\gamma_{q_j}(z)\alpha_j(a)$ for $a\in M_j$.
For a word $b_0[b_1|\cdots|b_n]$ homogeneous in
Euler weights $Q_0,\ldots,Q_n$, write
$S_i=Q_i+\cdots+Q_n$ and $D_Q=\prod_{i=1}^nS_i$.
The rational formula to be made integral is
\[
 \frac{\phi(b_0)\delta(b_1)\cdots\delta(b_n)}{D_Q}.
 \tag{AL.31}
\]
No division in $R$ is intended by this notation.
We now give its integral coefficients, which define
the map in every characteristic.

Expand each $b_i$ in the divided-power monomial basis.
For one such monomial, regard each occurrence of a
basis variable of Euler weight $q$ as a labelled
block of size $q$.  In entry $i$ let their total size
be $Q_i$, and put
$C_i=Q_i!/\prod q!$, the product running over those
blocks with multiplicity.  For $i>0$, each term of
the derivation selects one block of size $q_i$.
After multiplying the divided powers of $z$, the
coefficient of $\gamma_Q(z)$, $Q=\sum Q_i$, in
that term of \textup{(AL.31)} is
\[
 \frac{Q!\prod_{i=1}^nq_i}
      {D_Q\prod_{\text{all blocks}}q!}.
 \tag{AL.32}
\]
It is the following product of integers:
\[
 \binom{Q}{Q_0}C_0
 \prod_{i=1}^n
 \left[
   \binom{S_i-1}{Q_i-1}\frac{C_iq_i}{Q_i}
 \right].
 \tag{AL.33}
\]
Indeed
$C_iq_i/Q_i=(Q_i-1)!/
((q_i-1)!\prod_{\text{other blocks}}q!)$ is an
integer multinomial coefficient.  Multiplying
\textup{(AL.33)} telescopes the suffix factorials
and gives \textup{(AL.32)} exactly.
The terms remaining after this coefficient are
products of the specified $\alpha_j$'s and their
divided-power differentials, so are integral.
This defines \textup{(AL.31)} without division.
The same argument includes $n=0$, with an empty
product; a scalar word has coefficient one.

Work first in the free integral PD algebra on $z$
and on independent odd $x$'s with differentials the
corresponding even $y$'s.  It is torsion free.
The identities used to prove Theorem~\ref{AL05}
now hold with the positive Euler weights $Q_i$
in place of the arities $a_i$:
$d\delta(b)=Q_i\phi(b)$, merging two bar entries
adds their Euler weights, and every $\delta(b)$
is odd.  The coefficient telescope
$S_i-S_{i+1}=Q_i$ proves the chain identity, and
the labelled-shuffle identity with these weights
proves multiplicativity.  Both proofs are exactly
the scalar identities \textup{(AL.29)--(AL.30)}
with the indicated substitution.
The integral coefficients just proved and
injectivity into the rational algebra transfer
these identities and PD compatibility integrally.
Specialization then proves them in $R$.
On an arity-one generator in $M_j$, the formula
is the integral expression
$\gamma_{q_j}(z)\alpha_j$; the numerator factor
$q_j$ has already been cancelled as an integer
coefficient before any specialization.
All formulas commute with the announced parameter
and target maps, proving naturality.

For the balancing defect take
$k=\mathbb Z$, $S=\Gamma_{\mathbb Z}(T)$, and
\[
 R=S\otimes\Gamma_{\mathbb Z}(y_1,y_2)
          \otimes\Lambda_{\mathbb Z}(x_1,x_2),
 \qquad dx_i=y_i,\quad |x_i|=1,\ |y_i|=2.
\]
Let $Y_{q,i}$ denote a coefficient-indexed even
parameter whose image is $T^{[q]}y_i$ and whose
odd generator maps to $T^{[q]}x_i$.
Formula \textup{(AL.31)} gives
\[
 \begin{aligned}
 [Y_{1,1}Y_{1,2}]&\longmapsto
       T^{[2]}(x_1y_2+y_1x_2),\\
 [Y_{2,1}Y_{1,2}]&\longmapsto
       T^{[3]}(2x_1y_2+y_1x_2).
 \end{aligned}                                                   \tag{AL.34}
\]
If this map factored through balancing the parameters
over $S$, the module equation
$T Y_{1,1}=2Y_{2,1}$ would require the first image
multiplied by $T$ to equal twice the second.
Their difference is
\[
       T^{[3]}(-x_1y_2+y_1x_2)\ne0.                 \tag{AL.35}
\]
The two monomials have distinct members of the free
divided-power/exterior basis of $R$, so the
nonvanishing is proved over $\mathbb Z$.
For precision, this difference is the boundary
$d(T^{[3]}x_1x_2)$; our assertion is failure of
strict equality for this particular map, not
failure up to a specified homotopy, and not the
nonexistence of every possible balanced map.
\end{proof}
